\begin{table}[t]
\centering
\caption{Year-by-Year Bunching Estimates}
\label{tab:temporal}
\begin{tabular}{lcccc}
\hline\hline
 & \multicolumn{2}{c}{25-bed (all)} & \multicolumn{2}{c}{50-bed (non-CAH)} \\
\cmidrule(lr){2-3} \cmidrule(lr){4-5}
Year & $b$ & SE & $b$ & SE \\
\hline
2010 & 15.46 & (2.63) & 2.80 & (1.23) \\
2011 & 17.48 & (1.60) & 2.00 & (0.84) \\
2012 & 17.31 & (1.99) & 3.00 & (0.75) \\
2013 & 16.64 & (1.92) & 2.67 & (0.84) \\
2014 & 16.09 & (1.68) & 2.43 & (0.82) \\
2015 & 15.60 & (1.56) & 3.05 & (0.89) \\
2016 & 15.32 & (1.77) & 2.49 & (0.78) \\
2017 & 16.86 & (1.74) & 2.96 & (0.82) \\
2018 & 16.57 & (1.61) & 2.60 & (0.73) \\
2019 & 19.09 & (1.95) & 2.70 & (0.78) \\
2020 & 18.03 & (1.90) & 3.27 & (0.79) \\
2021 & 17.30 & (1.75) & 2.28 & (0.70) \\
2022 & 19.57 & (2.04) & 2.09 & (0.72) \\
2023 & 18.91 & (2.05) & 2.36 & (0.77) \\
\hline\hline
\end{tabular}
\par\vspace{0.3em}{\footnotesize \textit{Notes:} Annual bunching estimates at the 25-bed CAH threshold (all hospitals) and 50-bed RHC/REH threshold (non-CAH hospitals). Polynomial degree 7, manipulation window $[-2, +3]$ for 25-bed and $[-2, +3]$ for 50-bed. Standard errors from 100 bootstrap replications. The stability of $b$ across years confirms that bunching reflects persistent hospital sizing, not transient reporting noise.
}
\end{table}

